%% The following is a directive for TeXShop to indicate the main file
%%!TEX root = diss.tex

%=======================================
\chapter{The Algebraic Kernel Method}
\label{ch:AKM}
%=======================================

We will explore the algebraic kernel method as seen in \cite{bostan2014three} with extensions to interacting lattice walk models. We start with the non-interacting version of the problem.

%=======================================
\section{Models Without Interactions}
%=======================================
At the core of the Algebraic Kernel method we have the idea that we can apply group actions from $G$ onto the master functional equation and acquire nontrivially different equations for each $g\in  G$. After we sum these equations in a nice way, we manage to eliminate the functions $Q_1^0$, $Q_2^0$ and $\Qf{0,0}$ and their orbits, only leaving the orbit of the kernel term $KQ$ and a term on the right-hand side. What makes this possible is pre-multiplying the whole equation by $x_1 x_2$. This also works in higher dimensions as we will show later. Although the generalization has been done before (for example, page 11 of \cite{bousquet2010walks} for the two-dimensional case and page 10 of \cite{bostan2014three} for the three-dimensional one) we will present it in a compact and easy-to-generalize way. Instead of thinking just in terms of equations and summing them as in \cite{bousquet2010walks}, we will encapsulate the process in an operator $D_G$ that will take the functional equation and eliminate the terms we do not want by applying the group actions in a nice way. We start by defining $D_G$ so as to do the same thing as in \cite{bousquet2010walks} and \cite{bostan2014three} in order to formalize the process.
\begin{definition} Let $G$ be a group represented and generated by $g_k$'s preserving the kernel $K$ for a model for small steps. Then let $D_G$ be the operator defined as
\begin{align}
	D_G = \sum_{g\in G} (-1)^{|g|}g
\end{align}
in the sense of group actions. Here $|g|$ is the sign of the group element $g$.
\end{definition}
In two dimensions we can find the whole group by iterating over the two generators 1, $g_i$, $g_jg_i$, $g_ig_jg_i$. It is cleaner to write $g_{iji}$ instead of $g_ig_jg_i$, so we have the following notation.
\begin{definition} For a sequence $L = (l_1,l_2,\ldots)$ of positive integers, define 
\begin{align}
	g_L = g_{l_1l_2\ldots} = g_{l_1}g_{l_2}\cdots
\end{align} 
whenever this makes sense - when the list is finite or when there is a convergence in the case of an infinite list.
\end{definition}
In the case of a finite group in two dimensions we have $(g_1g_2)^n = 1$ for some minimal $n$. This allows us to rewrite $D_G$ in a useful way. The following lemma is just a restatement of the orbit sum method on page 11 of \cite{bousquet2010walks} but in operator notation.
\begin{lemma} For a two-dimensional model, with small steps and a group of order $2n$ with generators $g_1$ and $g_2$, we can write the operator $D_G$ as
\begin{align}
	D_G = (1+g_{ij}+\cdots+(g_{ij})^{n-1})(1-g_i)\,.
\end{align}
\end{lemma}
\begin{proof} We can just check the identity by expanding it as
\begin{align}\begin{split}
	D_G &= (1+g_{ij}+\cdots+(g_{ij})^{n-1})(1-g_i)\\
	        &= 1- g_i +g_{ij}-g_{iji} + \cdots - (g_{ij})^{n-1}g_i
\end{split}\end{align}
giving us one of the possible orbits
\end{proof}
Here are some examples of small groups.
\begin{example} When $|G| = 4$ the elements commute, so the orbit is $\{1$, $g_1$, $g_{12} = g_{21}$, $g_2\}$, and the operator is
\begin{align}
	D_G = 1 - g_1 - g_2 + g_{12} = (1-g_1)(1-g_2)\,.
\end{align}
When $|G| = 6$ the orbit is $\{1$, $g_1$, $g_{21}$, $g_{121}=g_{212}$, $g_{12}$, $g_2\}$. We can write the operator as
\begin{align}\begin{split}
	D_G &= (1+g_{12}+g_{1212})(1-g_1) \\
	&=  1+g_{12}+g_{1212} - g_1 - g_{121} - g_{12121} \\
	&= 1 -g_1 - g_2 + g_{12}+g_{21}-g_{121}\,.
\end{split}\end{align}
\end{example}
We can generalize the lemma for a multidimensional group $G$.
\begin{lemma}Let $\mc{M}$ be a lattice walk model with small steps in $\ze^d$, and let $G$ be its asociated group. For a subgroup $H$ of $G$ and a nontrivial generator $g_k$ of $G$, we have $D_G = D_H(1-g_k)$.
\end{lemma}
\begin{proof} Simpler proof: each element in the fixed representation of $G$ either ends in $g_k$ or does not. We have a bijection between the two sets by using the involution $g_k$ applied on the right. The elements ending in $g_k$ have opposite parity to those not ending on $g_k$, so we can split $G$ into two parts $H$ and $Hg_k$. We can split the defining sum into two parts and keep track of the sign changes via $e_g = (-1)^{|g|}$. We have
\begin{align}\begin{split}
	D_G &=\sum_{g \in H} e_g g + \sum_{g \in Hg_k} e_g g\\
	&=\sum_{g \in H} e_g g + \sum_{g \in H} (-1)e_g g g_k\\
	&= \left(\sum_{g \in H} e_g g\right)(1-g_k)\\
	&= D_H(1-g_k)\,.
\end{split}\end{align}
(PROOF) Take $H = G/\langle g_k \rangle$. Since $g_k$ is of order 2, the resulting group has $|G|/2$ elements. Consider the elements of $G$ that commute with $g_k$, that is, $h\in G$ where $g_k h = h g_k$. Multiplying by $g_k^{\um 1} = g_k$ we see that $h = g_k h g_k$. Let $H$ be the subgroup of $G$ generated by those elements. There are two cosets $H$ and $Hg_k$.
\end{proof}
\begin{lemma} If $I$ is a nonempty subset of $\{1,\ldots, d\}$, then the operator $D_G$ eliminates the term $t S_I^{-\one}Q_I^\oo$.
\end{lemma}
\begin{proof} By the previous lemma, for some $i\in I$,
\begin{align}
D_G(t S_I^{-\one}Q_I^\oo) = D_H((1-g_i)t S_I^{-\one}Q_I^\oo)) = D_H(0) = 0\,,
\end{align}
 since the variable $x_i$ is not present in the function.
\end{proof}
We acquire the $d$-dimensional generalization to Proposition 5 in \cite{bousquet2010walks}.
\begin{theorem} For a $d$-dimensional lattice walk model with small steps, we have
\begin{align}
D_G(\mon xd KQ) = D_G(\mon xd)\,,
\end{align}
where the right-hand side may be zero sometimes.
\end{theorem}
\begin{proof} When we multiply the master functional equation by $\mon xd$, we acquire a lot of terms like the ones in the lemma. The only term that survives the application of $D_G$ on the left-hand side is $\mon xd KQ$, and the equality follows.
\end{proof}

This is a generalization of the orbit summation found in \cite{melczer2014variants}. The authors in \cite{bousquet2010walks} and \cite{bostan2014three} have found that sometimes the algebraic kernel method fails. That is, $D_G(\mon xd) = 0$ and the equation does not give us new information. We can generalize this as follows.

\begin{corollary} Suppose that there is a group element $h \in G$ with odd parity that permutes the variables $\row xd$ and has order 2. Then the regular Algebraic Kernel method fails. That is,
\begin{align}
D_G(\mon xd KQ) = 0\,.
\end{align}
\end{corollary}
\begin{proof} Since $h$ is an involution, the same method applies as before and we can factor $D_G = D_H(1-h)$ for some $H \subseteq G$. The action $1-h$ eliminates the monomial $\mon xd$.
\end{proof}

\begin{example} The two-dimensional models for which the algebraic kernel method fails are precisely those where we have a permutation as in the corollary. That is, $h(x_1,x_2)= h(x_2,x_1)$ and $h$ has order three. Proposition 6 in \cite{bousquet2010walks} follows.
\end{example}

\begin{example} In three dimensions and models with at most 6 steps we see that the algebraic kernel method fails in 62 three-dimensional models (page 13 of \cite{bostan2014three}). For example, there is a model with group size 12 for which the algebraic kernel method fails. For that model, one of the group element is a permutation of two of the variables - in that case $x_2$ and $x_3$. Strangely enough, $G$ contain the group seen in the previous example for the variables $x_2$ and $x_3$. This means that every higher-dimensional model that has group $G$, containing the group above, will fail in the regular Algebraic Kernel method.
\end{example}
A natural question is to consider taking $D_H$ instead of $D_G = D_H(1-h)$ and checking if we can perfom a half-orbit summation. The answer is yes. The operator $D_H$ eliminates some of the $Q$ variables and introduces others, but the result is a function symmetric under permutation by $h$. One may hope to perform a method analogous to the half-orbit summation method to prove algebraicity or D-finiteness.

\begin{example}There is another way we can prove that $D_G$ does what it advertises. Like in \autoref{ch:HigherD}, we start with the kernel form $E(KQ) = 1$ where we make the substitution $Q = W/(\mon xd K)$ and rewrite the equation as
\begin{align}
	E' W = \mon xd\,, 
\end{align}
where $E' = E_1'\cdots E_d'$ and $E_k' = 1 - [x_k^0]$.
Therefore,
\begin{align}
	W = \mon xd -(1-E')W=\mon xd - \sum_{\emptyset \neq I \subseteq [d]}W_I^\oo\,.
\end{align}
We immediately see that
\begin{align}\begin{split}
	D_G(\mon xd KQ) = D_G(W) &= D_G\left(\mon xd - \sum_{\emptyset \neq I \subseteq [d]}W_I^\oo\right) \\&= D_G(\mon xd)\,,
\end{split}\end{align}
since the $W_I^\oo$ terms have missing variables. Let us think about this for a moment. The form of $W$ is the same for all lattice walk models with small steps so, if there are two models whose operators both annihilate expressions of the form $W_I^\oo$ for all nonempty $I\subseteq [d]$, then both models have equations of the above type with each other's operators! In particular, we may take the group $G$ of the highly-symmetric models and apply it via the operator $D_G$, for another model, and we retrieve an equation of the above type. We will look into this further.
\end{example}
\begin{definition} Let $g_k(x_k) = \ol x_k$. Define the operator $\Delta_k = 1 - g_k$ for this specific $g_k$. Let $\Delta = \mon \Delta d$.
\end{definition}

\begin{theorem} Let $G$ be the group generated by $g_k(x_k) = \ol x_k$ for $1 \leq k \leq d$. Then we have $D_G = \Delta$
and
\begin{align}
	\Delta(\mon xd KQ) = \Delta(\mon xd) = (x_1 - \ol x_1)\cdots(x_d-\ol x_d)
\end{align}
for all $d$-dimensional non-interacting lattice walk models with small steps in the positive orthant.
\end{theorem}
\begin{proof} It is enough to see that $\Delta$ annihilates terms of the form $W_I^\oo$ for $I \neq \emptyset$. Since the operators $\Delta_k$ commute, we only need to apply one $\Delta_i$ for some $i$ first, and find
\begin{align}
	\Delta_i(W_I^\oo) = W_I^\oo - g_i(W_I^\oo) = 0
\end{align}
since the term is missing variable $x_i$. The rest follows by linearity.
\end{proof}
This will not help to show a specific model has a D-finite solution because, in general, we cannot factor out the $K$ term from the left-hand side. However, the idea is important.

Next, we will look into generalizing the ideas involved in the algebraic kernel method.
\begin{definition} Let $G$ be a group acting on $\row xd$. Then we define a weighted orbit summation operator to be an operator of the form
\begin{align}
	D = \sum_{g\in G}R_g g\,,
\end{align}
where $R_g$ are functions in the variables $\row xd$.
\end{definition}

The orbit summation operator is an example of this with $R_g' = 1$ for the groups considered so far. A first view of this kind of operator can be seen on page 11 of \cite{wong2014confining} where the authors use an orbit summation with weights. We can capture various strategies in this notation
\begin{example} Since the Algebraic Kernel method fails to give a nice result when $G$ contains permutations, we may try to apply the operator $D_G$ on a slightly different equation - one where we multiply the whole equation by a function $T \neq \mon xd$. For a function $f$ we have
\begin{align}
	D_G(Tf) = \sum_{g\in G}g(T f) = \sum_{g\in G} ( g(T)) g(f) =  \sum_{g\in G} R_g g(f)\,.
\end{align}
The situation of applying $D_G$ on $T$ times the functional equation is the same as applying a weighted orbit summation operator on the original functional equation with weights $g(T)$ for $g \in G$.
\end{example}
Also, we may consider the strategy of applying a weighted operator several times in the hopes of attaining an easier equation to work with. However, as we shall see this is no different than looking at a weighted operator in the first place. This is because $G$ is a group and the composition $D_G^2$ leads to another weighted operator. Also, since some of the operators $D_G$ have nontrivial kernels they are not invertible, so the set of operators is not a group. However, we can show that the set is a monoid under composition.
\begin{prop} For a group $G$, the set of weighted orbit summation operators forms a monoid under composition.
\end{prop}
\begin{proof} Since $G$ is a group, we have an identity element 1, which can also be written as a weighted orbit summation operator with $R_g = 0$ for all $g\neq 1$ and $R_1 = 1$. Thus, the set of operator has an identity element. Since composition of operators is associative, we only need to check that the set of operators is closed under composition. We have
\begin{align}
	D_1D_2 = \sum_{g\in G}\sum_{h\in H} R_g gT_h h = \sum_{g\in G}\sum_{h\in H} R_g T_{gh} gh \sum_{g\in G}(\sum_{h\in H} R_{gh_{\um 1}} T_{g}) g\,,
\end{align}
which is also an operator of this form once we collect terms corresponding to the same group elements.
\end{proof}

We hope to apply a weighted operator $D$ to a problem in order to gain more insight into the structure of the generating function $Q$. Since all models with small steps and finite groups also have a regular orbit summation operator, we have two main interesting cases to investigate.
\begin{itemize}
\item[1.] When is $D_G(\mon xd) = 0$ and can we find weights so that we have a new operator $D$ that acts like $D_G$ but spares $\mon xd$?
\item[2.] Can we create a useful orbit summation operator when $G$ is infinite?
\end{itemize}
We can answer the first question in the negative when it comes to unmarked lattice walk models. In general, there is no other weighted orbital operator except the one with weights $R_g = 1$ that does not also eliminate $\mon xd$. The second case remain unclear.

Here is the strategy we will employ. If we apply the operator $D$ on the master functional equation, we need to eliminate all terms except $KQ$. That is, $D$ needs to annihilate all the $Q$ terms that are missing variables. It is enough to work on the fundamental functional equation for $W = \mon xd KQ$. To go back to the master functional equation case we recognize that $\mon xd$ just make it a weighted operator and we can remove the weight by pre-dividing by $\mon xd$. The resulting operator is also weighted.

\begin{lemma}A weighted operator $D$ exists such that $D( KQ) = D(1)$ if and only if a weighted operator $D'$ exists such that $D'(f_I^\oo)$ for generic functions $f_I^\oo$ which are missing the variables indexed by $I$, for all $I$.
\end{lemma}
\begin{proof} The discussion above provides the direction from left to right for $f = W$. To prove it for generic functions $f$ it suffices to note that $f_I^\oo$ and $f_J^\oo$ are independent from each other if $I \neq J$. That is, $D(f_I^\oo)$ and $D(f_J^\oo)$ do not share terms. The converse direction is proved by setting $f = W$ and setting $D = D'\mon{\ol x}{d}$, where $(D'\mon{\ol x}{d} )(f) = D'(\mon{\ol x}{d} f)$.
\end{proof}
This essentially turns it (?) into a linear algebra problem. We just need to apply $D$ on generic functions $W_I^\oo$, which are missing the variables indexed by $I$, collect the different results and make their coefficients zero. That is, for every $f = W_I^\oo$ we look at the orbit
\begin{align}
	\orb_G(f) = \{g(f)| g \in G\}
\end{align}
and collect coefficients in
\begin{align}
	0 = D(f) = \sum_{g\in G}R_g g(f) = \sum_{f' \in Orbit(f)} R_g' f'
\end{align}
where we set
\begin{align}
	R_g' = \sum_{\substack{g(f) = f'\\g\in G}} R_g = 0\,.
\end{align}
Thus, we reach a linear system of $|\orb_G(f)|$ equations in $|G|$ unknowns. Of course, we can do this for all $2^d - 1$ terms $f = W_I^\oo$ and hopefully find a system with a nontrivial solution. Here is a simple example
\begin{example} In $d = 2$ dimensions we have the generic functions $W(0,0) = const$, $W(0,x_2)$, $W(x_1,0)$ where at least one variable is missing. We can write them as $1, f(x_1), f(x_2)$. Now we consider the group generated by $g_k(x_k) = \ol x_k$
\begin{align}
	G = \{1, g_1, g_2, g_{12} = g_{21}\}\,.
\end{align}
We write the operator $D$ as 
\begin{align}
	D = R - R_1 g_1 - R_2 g_2 + R_{12} g_{12}\,.
\end{align}
Consider the different orbits. The one with the constant is trivial since the function does not change under the group action. Next we have
\begin{align}
	\orb_G(f(x_1)) = \{f(x_1),f(\ol x_1),f(\ol x_1), f(\ol x_1)\} = \{f(x_1),f(\ol x_1)\}
\end{align}
and similarly,
\begin{align}
	\orb_G(f(x_2)) = \{f(x_2),f(\ol x_2)\}\,.
\end{align}
This means we have the equations
\begin{align}
	R - R_1 &+ R_2 - R_{12} = 0\,,\\
	R = R_2,\! &\quad\, R_1 = R_{12}\,,\\
	R = R_1,\! &\quad\, R_2 = R_{12}\,.
\end{align}
Folowing the implications provided from these operators we find $R= R_1 = R_2 = R_{12}$. The simplest choice is to set $R = 1$ and we find that every weighted orbit summation operator for this group $G$ must be trivially-weighted if it were to annihilate generic functions with missing variables.
\end{example}
\begin{definition} If all the weights of a weighted operator are the same, we call it trivial or unweighted.
\end{definition}
\begin{lemma} Consider a $d$-dimensional non-interacting lattice walk model with small steps in the positive orthant. If its group is finite and has $d$ generators, then all orbital operators are trivial.
\end{lemma}
We need to check this proof.
\begin{proof}It is enough to show that $D(W_k^0) = 0$ for all $k$ implies that the operator is trivial. To this end suppose that $D(f) = 0$ where $f = W_k^0$. Then we can partition the group $G$ as before into sets $H$ and $Hg_k$. We see that $hg_k (f) = h(f)$, so the orbit of $f$ contains at most $|G|/2$ elements. However, these elements must all be distinct otherwise they would come from the same operations, but the group elements in $H$ are all distinct.Therefore, $|\orb(f) = |G|/2$. This means that we have the equations $R_h = R_{hg_k}$ for every $h \in H$. Once we go through every $k$ we have overlaps $R = R_1 =\cdots = R_d$ and $R_1 = R_{12} = R_{13}$, and so on. This means that all the $R$'s are the same and the operator is trivial.
\end{proof}
This shows that in general, there is no clever weighted summation operator that also satisfies the condition $D(\mon xd) \neq 0$. However, the idea will be fruitful in that we may find nontrivial partial solution in the interacting case.

%=======================================
\section{Models with interactions}
%=======================================
It is natural to consider extending the algebraic kernel method to one working for interacting models, but we quickly see that the regular orbit summation operator is not sufficient. This is where the power of the weighted orbital summation operators starts to show. The analysis for two-dimensional models is particularly nice because we end up having the same number of equations as unknowns, and we are able to find the required weights. The linear systems can be written as matrix equations where the determinant of the matrix tells us when we have a useful orbit summation operator. In the end the method does not extend fully because we do not eliminate the $\Qf{0,0}$ term. The D-finiteness of the full generating function $Q$ depends on the D-finiteness of the function $\Qf{0,0}$. Also, we hope to say something about the initial goal of establishing relationship between the full and partial interactions for a given model. We will only consider the interesting finite cases with 2 nontrivial generators. The group $G$ has order 4, 6 or 8, and the functional equation in two dimensions is
\begin{align}
	KQ  - K_1 Q_1^0 -  K_2 Q_2^0 + K_{1,2}Q_{1,2}^{0,0} = \alpha\,.
\end{align}
%====================================
\subsection{Groups of Order Four}
%====================================
From a previous section we found that
\begin{align}
	D_G = 1-g_1-g_2+g_1g_2 = (1-g_1)(1-g_2) = \Delta_1\Delta_2
\end{align}
where we defined $\Delta_i = 1- g_i$. In the non-interacting orbit summation method we had a direct application of $D_G$ on the equations to perform the required cancellations. Here we will attempt to do something similar for the interacting case. Since the cancellations do not happen automatically if we use $D_G$ directly, we will instead find a linear system of equations by applying the group operations on the original functional equation. Then we will use linear algebra to find a solution where our variables will be the ones we wish to eliminate $Q_i^0$ and their images under the group actions. If the associated matrix has a nonzero determinant, we have a unique solution and unfortunately, we will not be able to eliminate variables. However, there are some cases where the determinant is actually zero. Because of symmetry we expect to be able eliminate the desired variables. The group operations commute, so we have
\begin{align}
	g_i Q_i^0 = Q_i^0\quad\text{and}\quad g_ig_j Q_i^0 = g_j Q_i^0\,,
\end{align}
and therefore, we have the functions that we wish to eliminate
\begin{align}
	E = \{Q_1^0,\,Q_2^0,\,g_2Q_1^0,\,g_1Q_2^0\}\,.
\end{align}
This is looking promising since we have four equations and four unknowns that we wish to eliminate. We rewrite the original equation in a format centered on these variables
\begin{align}
	 K_1 Q_1^0 +  K_2 Q_2^0 = KQ + K_{1,2}Q_{1,2}^{0,0} -\alpha = \beta\,.
\end{align}
\begin{table}[H]
\caption{Effects of group actions $g\in G$ on $Q_1^0$ and $Q_2^0$.}
\label{tbl:actions4}
\centering
\begin{tabular}{ccc}
\toprule
	Group action & Effect on $Q_1^0$ & Effect on $Q_2^0$\\
\midrule
	$1$ & $Q_1^0$ & $Q_2^0$\\
	$g_1$ & $Q_1^0$ & $g_1Q_2^0$ \\
	$g_2$ & $g_2Q_1^0$ & $Q_2^0$ \\
	$g_{12}=g_{21}$ & $g_2Q_1^0$ & $g_1Q_2^0$ \\
\bottomrule
\end{tabular}
\end{table}
Therefore, we have the equations
\begin{align}
	 K_1 Q_1^0 +  K_2 Q_2^0 &= \beta\,,\\
	 g_1K_1 Q_1^0 + g_1K_2 g_1Q_2^0 &= g_1\beta = \beta_1\,,\\
	 g_2K_1 g_2Q_2^0 + g_2K_2 Q_2^0 &= g_2\beta = \beta_2\,,\\
	 g_{12}K_1 g_2Q_1^0 + g_{12}K_2 g_1Q_2^0 &= g_{12}\beta = \beta_{12}\,.
\end{align}
Using the order of variables in $E$ we can write the associated matrix $M$ of this linear system as
\begin{align}
	M_4 = \begin{bmatrix}
	K_1 & K_2 & 0 & 0 \\
	g_1 K_1 & 0 & 0 & g_1 K_2\\
	0 & g_2 K_2 & g_2 K_1 & 0\\
	0 & 0 & g_{12}K_1 & g_{12} K_2
	\end{bmatrix}
\end{align}
Its determinant is
\begin{align}\begin{split}
	\det M_4 &= K_1 g_{12}K_1 g_1K_2g_2K_2 -g_1K_1g_2K_1K_2g_{12}K_2\\
	& = K_1 g_{12}K_1 g_1K_2g_2K_2  -g_1(g_1K_1g_{12}K_1g_1K_2g_{2}K_2)\\
	&=\Delta_1  (K_1 g_{12}K_1 g_1K_2g_2K_2)\\
	&=\Delta_2 (g_1 K_1 g_2 K_1K_2 g_{12} K_2 )
\end{split}\end{align}
where the last equation comes from doing the process in the same way but with $g_2$ instead.

%====================================
\subsubsection{When is $\det M_4 = 0$?}
%====================================

Now our task is to find where this determinant is zero. Clearly the expression within each $\Delta$ must be in the kernel of the operator. To that end let us determine when that happens. For a Laurent polynomial $f$ in the variables $x_i$ and $t$, we have
\begin{align}
	\Delta_k f = \sum_n f_n (1-g_k)x_k^n = \sum_n (f_n - f_{- n}\mu_k^{- n})x_i^n
\end{align}
where $\mu_k = S_k^{\um 1}/S_k^1$. This implies that
\begin{align}
	f_{n} &= f_{- n} \mu_k^{- n}\\
	f_{n} (S_k^1)^n &= f_{- n}(S_k^{\um 1})^n
\end{align}
for all $n \in \ze$. It is not clear if this will help us. However, the models that exhibit groups of order 4 have at least one symmetry across a line - either vertical or horizontal symmetry. WLOG let the model have vertical symmetry. That is, 
\begin{align}
	g_1(x_1) = x_1, \quad g_1(S_2^i) = S_2^i,\quad g_1(S_1^i) = S_1^{-i}\,.
\end{align}
This allows us to compute some things. For example,
\begin{align}
	g_1 K_2 &= c_2 - t\ol x_2 g_1S_2^{\um 1} = c_2 - t\ol x_2 S_2^{\um 1} = K_2\,,
	g_1g_2K_2 = g_2 g_1 K_2 = g_2\,.
\end{align}
Therefore, $\Delta_1$ fixes both $K_2$ and $g_2K_2$.
We see that the determinant is
\begin{align}
	\det M_4 = \Delta_1  (K_1K_2 g_{12}K_1 g_2K_2) = K_2 g_2K_2 \Delta_1  (K_1 g_{12}K_1)\,.
\end{align}
This is zero if and only if $\Delta_1  (K_1 g_{12}K_1) = 0$, so we have
\begin{align}\begin{split}
	0&=\Delta_1  (K_1 g_{12}K_1) \\
	&= K_1g_{12}K_1 - g_1 K_1 g_2 K_2 \\
	&= (c_1-t\ol x_1S_1^{\um 1})(c_1-tx_1g_2S_1^{\um 1})-(c_1-tx_1S_1^{\um 1})(c_1-t \ol x_1 g_2S_1^{\um 1})\\
	&= c_1 t \ol x_1 (g_2 S_2^{\um 1} - S_2^{\um 1}) - c_1 t x_1 (g_2 S_2^{\um 1} - S_2^{\um 1})\\
	&= c_1 t (\ol x_1-x_1)(g_2 S_1^{\um 1}-S_1^{\um 1})\\
	& = -c_1 t(\ol x_1-x_1)\Delta_2(S_1^{\um 1})\,.
\end{split}\end{align}
This is true if and only if
\begin{align}
	g_2 S_1^{\um 1} = S_1^{\um 1}\,.
\end{align}
We have
\begin{align}
	g_2 (\Sf{\um 1, \um 1}\ol x_2+\Sf{\um 1,0}+\Sf{\um 1,1}x_2) = \Sf{\um 1, \um 1}x_2 \ol \mu_2+\Sf{\um 1,0}+\Sf{\um 1,1}\ol x_2\mu_2
\end{align}
and
\begin{align}
	\Sf{\um 1, \um 1}x_2 \ol \mu_2+\Sf{\um 1,1}\ol x_2\mu_2 = \Sf{\um 1, \um 1}\ol x_2+\Sf{\um 1,1}x_2
\end{align}
that can be re-written succinctly as
\begin{align}
	\Sf{\um 1, \um 1} = \mu_2 S_{\um 1,1}\,.
\end{align}
If one of the $S$'s is zero, so is the other, and therefore, by symmetry, the model does not have diagonal steps. Since we have two nontrivial generators, this implies that both $x_2$ and $\ol x_2$ are present in $S$, so the model trivially exhibits a symmetry about the $x_1$ axis. If the $S$'s are nonzero, this implies $\mu_2$ is a constant and therefore, $S_2^{\um 1} = S_2^1$. Again the model exhibits symmetry about the $x_1$ axis. 
\begin{theorem} We have that $\det M_4 = 0$ if and only if the model is highly-symmetric.
\end{theorem}

\begin{corollary} The algebraic kernel method extends to the non-interacting case for $|G| = 4$ if and only if the model is highly-symmetric.
\end{corollary}
\begin{figure}[H]
    \centering
\[
	\modelgraphB 01011010 \quad \modelgraphB10100101  \quad \modelgraphB 11100111  \quad \modelgraphB 10111101  \quad \modelgraphB 11111111 
\]
    \caption{The highly-symmetric models on the plane.}
\end{figure}
We are now ready to attempt using the orbit summation method in the highly symmetric cases. There are 8 such models, but only 5 of them have dimension 2. Furthermore, 2 of the models are reflections of each other, so we are really left with 4 interesting cases - diagonal walks, full walks, full walks minus two horizontal steps. If we are able to eliminate all the variables, we will be left with a linear relation like
\begin{align}
	\gamma \beta- \gamma_1 \beta_1-\gamma_2 \beta_2 + \gamma_{12}\beta_{12} = 0\,.
\end{align}
Here the choice of signs is inspired by the orbit summation operators. Let $\bullet$ and $\ast$ stand for applications of $g_1$ and $g_2$, respectively. First we have $K_2^\bullet = K_2$ and $K_1^\ast = K_1$. Then we let
\begin{align}
	x = Q_1^0\,,\quad y= Q_2^0\,, \quad u= Q_1^{0,\ast}\,,\quad v = Q_2^{0,\bullet}
\end{align}
and
\begin{align}
	a = K_1\,,\quad b = K_1^\bullet\,,\quad c=K_2 \,, \quad d = K_2^\ast
\end{align}
The original equations become
\begin{align}
	\beta &= a x + c y\,,\quad \beta_1 = b x + c v\,,\\
	\beta_2 &= d y + a u\,,\quad \beta_{12} = b u + d v\,,
\end{align}
Therefore, the above equation needs to satisfy
\begin{align}\begin{split}
	0 &= (ax+cy)\gamma - (bx+cv)\gamma_1 - (dy+au)\gamma_2 + (bu +dv) \gamma_{12}\\
	   &= (a\gamma-b\gamma_1)x + (c\gamma-d\gamma_2)y + (b\gamma_{12}-a\gamma_2)u + (d\gamma_{12}-c\gamma_1)v
\end{split}\end{align}
and in order to eliminate the variables $x,y,u,v$ we need to make all their coefficients zero. Therefore,
\begin{align}
	a\gamma = b\gamma_1\,,\quad c\gamma = d\gamma_2\,, \quad b\gamma_{12} = a \gamma_2,\,,\quad d\gamma_{12} = c\gamma_1\,.
\end{align}
If we set $\gamma = 1$, then $\gamma_1 = a/b$, $\gamma_2 = c/d$ and $\gamma_{12} = (ac)/(bd)$. So we can multiply out by $bd$ to find the relation
\begin{align}
	bd\beta-ad\beta_1-cb \beta_2+ac \beta_{12} = 0
\end{align}
where we have successfully eliminated $Q_1^0, Q_2^0, Q_1^{0,\ast}$, and $  Q_2^{0,\bullet}$. Or written in terms of the $K$'s we have
\begin{align}
	K_1^{\bullet}K_2^{\ast}\beta^{}-K_1^{}K_2^{\ast}\beta^{\bullet}-K_1^{\bullet}K_2^{}\beta^{\ast}+K_1^{}K_2^{}\beta^{\bullet\ast} = 0\,.
\end{align}
We can rewrite this equation in the following nicer way.
\begin{theorem} For highly-symmetric models with small steps, we have the functional equation
\begin{align}
	\Delta(K_1^\bullet K_2^\ast \beta) = 0\,,
\end{align}
where $\beta = KQ + K_{1,2}Q_{1,2}^{0,0} -\alpha$.
\end{theorem}
This is a generalization of the result in \cite{wong2014confining} where the authors use a weighted orbit sum to eliminate boundary terms. We can expand the functional equation to produce a different form. Towards this end let $K_0 = K_1^\bullet K_2^\ast$. We find
\begin{align}\begin{split}
	\Delta(K_0\beta) &= \Delta(K_0(KQ + K_{1,2}Q_{1,2}^{0,0} - \alpha)) \\
	&= \Delta(K_0 KQ) +\Delta(K_0 K_{1,2}Q_{1,2}^{0,0}) - \Delta(K_0\alpha)\\
	&= K\Delta(K_0 Q) + \Delta(K_0 K_{1,2})Q_{1,2}^{0,0} - \Delta(K_0)\alpha
\end{split}\end{align}
where the last line follows since $\Delta$ commutes with $K$ and functions with some variables $x_i$ missing. All that remains is to expand the various quantities involved
\begin{align}
	\Delta g_iK_i = - \Delta K_i = - (c_i-tx_iS_i^{\um 1}-c_i + t\ol x_i S_i^{\um 1}) = -t(\ol x_i-x_i)S_i^{\um 1}\,.
\end{align}
Therefore,
\begin{align}
	\Delta K_0 = \Delta_1 (K_1^\bullet) \Delta_2 (K_2^\ast) = t^2(\ol x_1-x_1)(\ol x_2-x_2)S_1^{\um 1}S_2^{\um 1}\,.
\end{align}
Next
\begin{align}
	\Delta(K_0 K_{1,2}) = \Delta_1 K_1^\bullet \Delta_2 K_2^\ast K_{1,2} = \Delta_1 K_1^\bullet (K_2^\ast K_{1,2}- K_2 K_{1,2}^\ast)\,.
\end{align}
The quantity in parentheses is
\begin{align}\begin{split}
	&K_2^\ast K_{1,2}- K_2 K_{1,2}^\ast =\\
	=& (c_2-tx_2S_2^{\um 1})(c_1c_2-t\ol x_1\ol x_2S_{1,2}^{\um 1, \um 1})-(c_2-t\ol x_2S_2^{\um 1})(c_1c_2-t\ol x_1x_2S_{1,2}^{\um 1, \um 1})\\
	=& c_1c_2tS_2^{\um 1}(\ol x_2-x_2) - c_2t\ol x_2S_{1,2}^{\um 1, \um 1}(\ol x_2-x_2)\\
	=& c_2t(\ol x_2-x_2)(c_1S_2^{\um 1}-\ol x_1S_{1,2}^{\um 1, \um 1})\,.
\end{split}\end{align}
and thus,
\begin{align}\begin{split}
	\Delta(K_0 K_{1,2})  &= c_2t(\ol x_2-x_2) \Delta_1 K_1^\bullet (c_1S_2^{\um 1}-\ol x_1S_{1,2}^{\um 1, \um 1})\\
	&=c_2t(\ol x_2-x_2) \Delta_1 K_1^\bullet (c_1S_2^{\um 1}-\ol x_1S_{1,2}^{\um 1, \um 1})\,.
\end{split}\end{align}
To simplify this quantity note that
\begin{align}
\Delta_1 K_1^\bullet c_1 S_2^{\um 1} = c_1 S_2^{\um 1}\Delta_1 K_1^\bullet = -c_1t S_2^{\um 1}(\ol x_1-x_1)
\end{align}
and
\begin{align}
	\Delta_1 K_1^\bullet \ol x_1S_{1,2}^{\um 1, \um 1} = S_{1,2}^{\um 1, \um 1}\Delta_1 K_1^\bullet \ol x_1 = c_1 (\ol x_1-x_1)S_{1,2}^{\um 1, \um 1}\,.
\end{align}
Finally,
\begin{align}
	\Delta(K_0 K_{1,2}) = c_1c_2t(\ol x_1-x_1)(\ol x_2-x_2)(t S_1^{\um 1} S_2^{\um 2}-S_{1,2}^{\um 1, \um 1})\,.
\end{align}
Combining terms and simplifying a bit leads us to
\begin{align}
	K\Delta(K_0 Q) = \alpha t(\ol x_1-x_1)(\ol x_2-x_2)(tS_1^{\um 1}S_2^{\um 1} - b_1b_2(tS_1^{\um 1}S_2^{\um 1}-\Sf{\um 1, \um 1})\Qf{0,0})
\end{align}
\begin{corollary} For a highly-symmetric model with $K_0 = K_1^\bullet K_2^\ast$, we have the functional equation
\begin{align}
	K\Delta(K_0 Q) = \alpha t(\ol x_1-x_1)(\ol x_2-x_2)(tS_1^{\um 1}S_2^{\um 1}(1 - b_1b_2\Qf{0,0}) + \Sf{\um 1, \um 1}\Qf{0,0})\,.
\end{align}
\end{corollary}
%====================================
\subsection{Groups of Order Six}
%====================================
Since we have 6 group actions, it is worth taking the time to see what kinds of expressions we develop by applying $g$ on $Q_i^0$ or $K_i$ for each $i$. For example, $g_i Q_i^0 = Q_i^0$, so we say the action of $g_i$ is reduced to $g_{\emptyset} = 1$ on $Q_i^0$. Using the sequences $L$ produced in one orbit of a representation of $G$ and their alternatives, we can tabulate their reductions on the different expresions. For example,
\begin{align}
g_{2121}Q_1^0 = g_{212}Q_1^0 = g_{121} Q_1^0 = g_{12}Q_1^0
\end{align}
where on the second step we have used the alternative form $g_{212} = g_{121}$. The table has $12$  in the row and column corresponding to $2121$ and $Q_1^0$. We omit the row corresponding to the identity.
\begin{table}[H]
    \caption{Reductions of group actions $g \in G$ on $Q_i^0$ and $K_i$.}
    \label{tbl:tableG6}
    \centering
    \begin{tabular}{cccccc}
    \toprule
    Sequence $L$ & Alternative & $Q_1^0$ & $Q_2^0$ & $K_1$ & $K_2$ \\ 
    \midrule
    1 & 21212 & $\emptyset$ & 1& 1& 1 \\
    21 &  1212& 2& 21 & 21& 21  \\
    121 & 212 & 12& 21& 121& 212 \\
    2121 & 12 &  12 & 1 & 12 & 12 \\
    12121 & 2 &  2& $\emptyset$& 2& 2 \\
    \bottomrule
    \end{tabular}
\end{table}
We only have 6 different functions results of applying $G$ on $Q_1^0$ and $Q_2^0$, so we have only 6 variables that can be eliminated.  Let
\begin{align}
	u_1 &= Q_1^0\,,\quad u_2 = g_2 Q_1^0\,,\quad u_3 = g_{12}Q_1^0\\
	v_1 &= Q_2^0\,,\quad v_2 = g_1 Q_2^0\,,\quad v_3 = g_{21}Q_2^0\,.
\end{align}
Also, there are 12 different $K_1$'s and $K_2$'s that we label sequentially as $\row A6$ and $\row B6$ respectively. The associated matrix to the linear system of equations is
\begin{align}
	M_6 = \begin{bmatrix}
	A_1& & & B_1 & & \\
	A_2& & & & B_2& \\
	& A_3& & & & B_3\\
	& & A_4 & & & B_4\\
	& & A_5& & B_5& \\
	&A_6 & &B_6 & & \\
	\end{bmatrix}
\end{align}
and it has determinant
\begin{align}
	\det M_6 = A_2A_4A_6 B_1 B_3 B_5 - A_1 A_3 A_5 B_2 B_4 B_6\,.
\end{align}
The even index $A's$ are found by applying $g_1$ on the odd index $A's$ and vice versa for the $B's$. That is,
\begin{align}
	\det M_6 = \Delta_1 (A_2A_4A_6 B_1 B_3 B_5) = \Delta_1(g_1K_1 g_{121}K_1g_{12}K_1 K_2 g_{21}K_2 g_{12}K_2)\,.
\end{align}
In the following table we show the relations between $A$'s and $B$'s for the 6 models with one of the two groups of order 6. $S11$, $S12$ and $S13$ have the one group; and $S21$, $S22$ and $S23$, have the other.
\begin{table}[H]
\caption{$A = K_1$ and $B = K_2$ relations for groups of order 6.}
\label{tbl:Relations6}
\centering
\begin{tabular}{cccc}
\toprule
	Model & Name & $A$-relations & $B$-relations \\ 
\midrule
	\centering \modelgraph 01010001 & S11 & $A_1 = A_6$, $A_2 = A_5$, $A_3 = A_4$ & $B_1 = B_4$, $B_2 = B_3$, $B_5 = B_6$\\
	\centering \modelgraph 10001010 & S12 &  $A_1 = A_4$, $A_2 = A_3$, $A_5 = A_6$ & $B_1 = B_2$, $B_3 = B_6$, $B_4 = B_5$\\
	\centering \modelgraph 11011011 & S13& None & None\\
	\centering \modelgraph 00110010 & S21  &  $A_1 = A_6$, $A_2 = A_5$, $A_3 = A_4$ & $B_1 = B_2$, $B_3 = B_6$, $B_4 = B_5$\\
	\centering \modelgraph 01001100 & S22 & $A_1 = A_4$, $A_2 = A_3$, $A_5 = A_6$ & $B_1 = B_4$, $B_2 = B_3$, $B_5 = B_6$\\
	\centering \modelgraph 01111110 & S23 & None & None\\
    \bottomrule
   \end{tabular}
\end{table}

When a model has nontrivial relations, then $\det M_6 = 0$ and, using a computer, we can easily check that $\det M_6 \neq 0$ when there are no relations.

We seek the same type of relation on $\beta$ and its orbit as for groups of order 4, so we will use $D_G$ as inspiration for the sign changes. We have the following
\begin{align}\begin{split}
0=&\gamma\beta-\gamma_{1}\beta_{1}+\gamma_{21}\beta_{21}-\gamma_{121}\beta_{121}+\gamma_{12}\beta_{12}-\gamma_{2}\beta_{2}\\
	=&\gamma (A_1 u_1 + B_1 v_1) - \gamma_{2} (A_6 u_2 + B_6 v_1) - \gamma_{1} A_2 u_1 + B_2 v_2) \\&+  \gamma_{12} (A_5 u_3 + B_5 v_2) + \gamma_{21}(A_3 u_2 + B_3 v_3) - \gamma_{121} (A_4 u_3 + B_4 v_3)\,.
\end{split}\end{align}
Expanding out and collecting terms produces the system
\begin{align}
	 A_1 \gamma&= A_2 \gamma_{1}\,,\quad   \quad A_3 \gamma_{21} = A_6 \gamma_{2}\,,\quad \quad A_5 \gamma_{12} = A_4 \gamma_{121}\,,\\
	 B_1 \gamma &= B_6\gamma_{2}\,,\quad  \quad  B_5 \gamma_{12}= B_2 \gamma_{1}\,,\quad  \quad B_3 \gamma_{21} = B_4 \gamma_{121}\,.
\end{align}
This system of equations does not have a nontrivial symbolic solution, so we need to use the relations here. There is a rational solution in each case with relations where $\gamma$ is a free variable. We set $\gamma$ to be the least common multiple of the denominators. We find that $K_0 = \gamma$ in each of the four cases with relations.

\begin{table}[H]
    \caption{Weight $K_0$ for the Orbit Summation operators.}
    \label{tbl:tableWeight6}
    \centering
   \begin{tabular}{cccc}
    \toprule
    Model & $S$-polynomial & Weight $\gamma = K_0$ & $K$-form \\ 
    \midrule
	\centering \modelgraph 01010001 & $ x_1\ol x_2+\ol x_1+x_2$  & $A_3 A_5 B_6$ & $ g_{21}(K_1) g_{12}(K_1) g_{2}(K_2)$\\
	\centering \modelgraph 10001010 & $\ol x_1x_2 + x_1 + \ol x_2$ & $A_2 A_5 B_6$ & $ g_{1}(K_1) g_{12}(K_1) g_{2}(K_2)$\\
	\centering \modelgraph 00110010 & $ x_1 x_2+\ol x_1+\ol x_2$ & $A_2 A_3 B_5 B_6$ & $g_{1}(K_1) g_{21}(K_1) g_{12}(K_2) g_{2}(K_2)$\\
	\centering \modelgraph 01001100 & $\ol x_1\ol x_2 + x_1 + x_2$ & $A_2 B_6$ & $g_{1}(K_1) g_{2}(K_2)$\\
    \bottomrule
   \end{tabular}
\end{table}
Here we have picked the respresentations of $A$'s and $B$' with the least number of group operations $g_i$ to reach the resulting expression. We also tried finding a possible $K_0$ by brute force for $S13$ and $S23$, but we did not find any among products of distinct $A$'s and $B$'s. It is likely that such a weight does not exist for these cases.
\begin{theorem} The algebraic kernel method extends for interacting lattice walk models with small steps and group $G$ of size 6 if and only if $\det M_6 = 0$. In those cases we have the functional equation
\begin{align}
	D_G(K_0 \beta) = 0\,,
\end{align}
where $\beta =  KQ + K_{1,2}Q_{1,2}^{0,0} -\alpha$ and $K_0$ depends on the specific model.
\end{theorem}
%====================================
\subsection{Groups of Order Eight}
%====================================
The analysis for groups of order 8 is quite similar, so we omit it. Here are the results. 
\begin{table}[H]
    \caption{Models with group of size 8.}
    \label{tbl:tableWeight6}
    \centering
   \begin{tabular}{cc}
    \toprule
    Model & $S$-polynomial \\ 
    \midrule
	\centering \modelgraph 10011001 & $x_1+\ol x_1+x_1\ol x_2+\ol x_1x_2$\\
	\centering \modelgraph 00111100 & $x_1+\ol x_1+x_1x_2+\ol x_1\ol x_2$\\
    \bottomrule
   \end{tabular}
\end{table}

We name the $A_1 = K_1$ and $B_1 = K_2$ iterates as before; for example, \ $A_2 = g_1 A_1$ and $A_3 = g_2 g_1 A_1$.
\begin{table}[H]
\caption{Reductions of group actions $g \in G$ on $Q_i^0$ and $K_i$.}
\label{tbl:tableG6}
\centering
\begin{tabular}{cccc}
\toprule
	Sequence $L$ & Alternative & $Q_1^0$ & $Q_2^0$ \\ 
\midrule
	1 & 2121212 & $\emptyset$ & 1 \\
	21 &  121212& 2 & 21   \\
	121 & 21212 & 12 & 121  \\
	2121 & 1212 &  212 & 121  \\
	12121 & 212 &  212 & 21  \\
	212121 & 12 &  12& 1  \\
	1212121 & 2 &  2& $\emptyset$ \\
\bottomrule
\end{tabular}
\end{table}
The relations are the same in both models
\begin{align}
	A_1 &= A_6\,,\quad A_2 = A_5\,,\quad A_3 = A_4\,,\quad A_7 = A_8\,\\
	B_1 &= B_4\,, \quad B_2 = B_3\,,\quad B_4 = B_8\,,\quad B_6 = B_7\,.
\end{align}
In both cases we retrieved the same $\gamma$ operator that we have written in a form involving the least number of group actions.
\begin{align}
	\gamma = K_0 =  A_2 A_3 B_5 B_6 = g_1(K_1) g_{21}(K_1) g_{212}(K_2) g_{1212}(K_2)\,.
\end{align}
\begin{theorem} The algebraic kernel method extends for all interacting lattice walk models with small steps and group $G$ of size 8. The functional equation is
\begin{align}
	D_G(K_0 \beta) = 0\,,
\end{align}
where $\beta =  KQ + K_{1,2}Q_{1,2}^{0,0} -\alpha$ and $K_0 = g_1(K_1) g_{21}(K_1) g_{212}(K_2) g_{1212}(K_2)$.
\end{theorem}

\endinput